% preamble-english.tex
% -----------------------------------------------------------------------
% Lo propio de "Read and Draw" (english.tex, english-bw.tex), el cuaderno
% en inglés -- ver notes/05-english.md. Se carga DESPUÉS de preamble.tex
% (con \booklang = en: babel en inglés británico y las cadenas de
% lang/en.tex) y carga a su vez preamble-lupa.tex: este cuaderno es un
% nivel por debajo de "Leo con lupa" y comparte con él la página entera
% -- la cabecera, el pie "Day N / 260", la caja de actividad que llena
% lo que queda de página, el texto en párrafos con sus notas y listas.
% Aquí solo lo que cambia:
%
%   - el tamaño del texto del día (\fuenteIngles): textos más cortos que
%     en "Leo con lupa", en un idioma que todavía se está aprendiendo --
%     letra más grande, de 20 a 17 pt;
%   - la caja de actividad (actividadIngles): la misma que actividadLupa,
%     con la letra de dentro más grande (15 pt, lo que lee la niña);
%   - las piezas de las actividades de tools/english.py: sí o no, rodear,
%     unir, los huecos con su caja de palabras, los dibujos para colorear
%     y las escenas donde dibujar (diagrams/english/);
%   - un icono por actividad, delante del título de su caja y en la
%     página "The words in this book" (frontmatter/english/palabras.tex):
%     quien empieza a leer en inglés reconoce antes un lápiz que la
%     palabra "draw".
% -----------------------------------------------------------------------
% preamble-lupa.tex
% -----------------------------------------------------------------------
% Lo propio del nivel 3, "Leo con lupa" (lupa.tex, lupa-bw.tex), cargado
% DESPUÉS de preamble.tex -- igual que preamble-palabras.tex para el
% nivel 1: la letra, la paleta, las cajas base, \cabeceraDia y la caja
% de lectura son las mismas que en el cuaderno de frases; aquí solo lo
% que cambia en este nivel (ver notes/04-nivel-lupa.md):
%
%   - la página: la caja de actividad llena todo lo que queda de página
%     (tcolorbox `height fill`) -- el sitio para dibujar, pensar o
%     escribir se adapta solo a lo largo que sea el texto de ese día.
%     Como en el cuaderno de primeras palabras, eso obliga a sacar el
%     pie "Día N / 260" del texto y ponerlo en el pie de página de
%     verdad (estilo `pielupa`);
%   - el tamaño del texto del día (\fuenteLupa): párrafos de 55 a 125
%     palabras, nunca por debajo de 14 pt;
%   - las piezas de las actividades de análisis, que genera tools/lupa.py
%     (dibuja con lupa, mapa, tabla de pistas, caso, errores, mensaje
%     secreto, ficha, compara, viñetas...);
%   - la lupa del club, en TikZ.
% -----------------------------------------------------------------------

% El mismo margen inferior que el cuaderno de primeras palabras
% (preamble-palabras.tex): con el pie a 10 mm del texto, la línea del
% pie queda a ~12 mm del borde del papel, dentro de lo que imprime
% cualquier impresora doméstica.
\geometry{bottom=22mm, footskip=10mm}

\makeatletter
\def\ps@pielupa{%
  \let\@oddhead\@empty\let\@evenhead\@empty
  \def\@oddfoot{\hfil\footnotesize\color{colorGris}\lblDia~\numeroDia{\lupapagina@actual}~/ \totaldias\hfil}%
  \let\@evenfoot\@oddfoot
}

% La página de un día en este cuaderno: misma cabecera que los otros dos
% (\cabeceraDia, preamble.tex, que pone también \label{dia:N} -- así
% tools/check_pages.py funciona igual sobre el .aux de los tres). \gdef
% y no \def: el pie se compone en la rutina de salida, y así no depende
% de que esa rutina corra todavía dentro del grupo del entorno.
\newenvironment{lupapagina}[4]{%
  \gdef\lupapagina@actual{#1}%
  \thispagestyle{pielupa}%
  \cabeceraDia{#1}{#2}{#3}{#4}%
}{%
  \par\newpage
}
\makeatother

% Tamaño del texto del día, según el trimestre (tools/libros.py, LUPA).
\newcommand{\fuenteLupa}[1]{%
  \ifcase#1\relax
  \or \fontsize{16}{21.5}\selectfont
  \or \fontsize{15}{20.5}\selectfont
  \or \fontsize{14.5}{19.5}\selectfont
  \or \fontsize{14}{19}\selectfont
  \fi
}

% Tamaño de letra de todo lo que la niña lee dentro de una actividad:
% más pequeño que el texto del día, pero nunca "de adulto" -- las
% instrucciones para el adulto van aparte, en pequeño y gris.
\newcommand{\letraActividad}{\fontsize{13}{17}\selectfont}

% Una sola caja para todas las actividades (color y título como
% argumentos). height fill: la caja llega hasta el final de la página.
% Lo que va tras \tcblower (la lista para comprobar un dibujo, las
% líneas de escribir, el banner de las páginas leídas) queda pegado al
% fondo: todo el espacio sobrante se lo lleva la parte de arriba (space
% to upper), sin línea separadora. Una caja de altura fija no respeta un
% \vfill dentro, por eso se hace así (ver preamble-palabras.tex).
%
% height plus: sin él, cuando el contenido NO cabe en el alto que le
% queda a la página, tcolorbox mantiene la altura fija y dibuja la
% parte de abajo ENCIMA de la de arriba -- sin ningún aviso en el log,
% ni Overfull ni nada: visto en el día 37 al bajar el margen inferior a
% 22 mm. Con height plus, la caja crece hasta su altura natural, no
% cabe en la página, y el día pasa entero a la siguiente -- que es justo
% lo que tools/check_pages.py detecta. Un desbordamiento visible y
% comprobado en vez de uno silencioso.
\newtcolorbox{actividadLupa}[2]{cajabase,
  colback=#1Claro, colframe=#1, colbacktitle=#1, title={#2},
  height fill, height plus=100cm, space to upper, segmentation hidden, middle=1mm,
  fontupper=\letraActividad, fontlower=\letraActividad,
  before upper={\hyphenpenalty=10000\exhyphenpenalty=10000},
  before lower={\hyphenpenalty=10000\exhyphenpenalty=10000}}

% --- dentro del texto del día ---------------------------------------
\newcommand{\separaParrafo}{\par\vspace{0.4em}}
% Una nota, un cartel o una carta que los personajes leen dentro de la
% historia ("> ..." en el JSON): recuadro blanco, para que se vea que es
% otro texto dentro del texto.
\newtcolorbox{notaTextoCaja}{enhanced, colback=white, colframe=colorGris,
  boxrule=0.6pt, arc=1.5mm, left=3mm, right=3mm, top=1.5mm, bottom=1.5mm,
  before skip=1mm, after skip=1mm}
\newcommand{\notaTexto}[1]{\begin{notaTextoCaja}\raggedright #1\end{notaTextoCaja}}
\newenvironment{listaTexto}{%
  \begin{itemize}[leftmargin=7mm, itemsep=0mm, topsep=0.2em, label=\textbullet]%
}{\end{itemize}}

% --- piezas de las actividades --------------------------------------
\newcommand{\renglon}{\par\vspace{8.5mm}\noindent\rule{\linewidth}{0.4pt}}
\newcommand{\renglonCorto}{\rule{0.6\linewidth}{0.4pt}}
\newcommand{\renglonNumerado}[1]{\par\vspace{8.5mm}\noindent\makebox[6mm][l]{#1.}\rule{\dimexpr\linewidth-6mm\relax}{0.4pt}}
\newcommand{\renglonCorrige}{\par\vspace{7mm}\noindent\hspace{6mm}\rule{\dimexpr\linewidth-6mm\relax}{0.4pt}}
\newcommand{\bannerDia}[1]{\par\vspace{2mm}{\centering\bfseries\Large\color{colorCrea}#1\par}}
\newcommand{\casillaMarca}{\raisebox{-0.3ex}{\framebox[4.6mm][c]{\rule{0pt}{3.6mm}}}}
\newcommand{\preguntaNumerada}[2]{\par\noindent\hangindent=6mm\hangafter=1\makebox[6mm][l]{\textbf{#1.}}#2\par}
\newcommand{\afirmacionVF}[1]{\par\vspace{2.5mm}\noindent
  \begin{tabularx}{\linewidth}{@{}X r@{}}#1 & \lblVerdadero\ \casillaMarca\quad\lblFalso\ \casillaMarca\end{tabularx}}
\newenvironment{listaNumerada}{%
  \begin{enumerate}[leftmargin=7mm, itemsep=0.5mm, topsep=1mm, label=\textbf{\arabic*.}]%
}{\end{enumerate}}
\newenvironment{listaOrdena}{%
  \begin{itemize}[label=, leftmargin=2mm, itemsep=5mm, topsep=3mm]%
}{\end{itemize}}

% Caso: las opciones van grandes y separadas, para rodear la buena.
\newcommand{\opcionCaso}[1]{\mbox{\Large\bfseries #1}}

% Errores: el resumen equivocado, en blanco y con interlineado amplio
% para poder subrayar o rodear encima.
\newtcolorbox{cajaRelato}{enhanced, colback=white, colframe=colorResponde,
  boxrule=0.6pt, arc=1.5mm, left=3mm, right=3mm, top=2mm, bottom=2mm,
  fontupper=\fontsize{14}{22}\selectfont\raggedright}
\newcommand{\relatoErrores}[1]{\begin{cajaRelato}#1\end{cajaRelato}}

% Tabla lógica: casillas grandes para escribir ✓ o ✗.
\newcommand{\cabeceraLogica}[1]{\rule{0pt}{5.5mm}\bfseries #1}
\newcommand{\celdaLogica}{\rule[-4mm]{0pt}{11mm}}

% Mapa en cuadrícula (filas con número, columnas con letra, como en
% "hundir la flota"). Las casillas se calculan en tools/lupa.py.
\tikzset{
  lineaMapa/.style={colorGris, line width=0.6pt},
  etiquetaMapa/.style={font=\bfseries\large, text=colorGris},
  lugarMapa/.style={font=\fontsize{11}{12.5}\selectfont, align=center},
}
\newcommand{\salidaMapa}{\textcolor{colorRelaciona}{\bfseries\lblSalida}}

% Mensaje secreto.
\newenvironment{tablaCodigo}{%
  \normalsize\setlength{\tabcolsep}{1.9mm}\renewcommand{\arraystretch}{1.3}%
  \begin{tabular}{|*{14}{c|}}%
}{\end{tabular}}
\newcommand{\celdaCodigo}[1]{%
  \begin{tabular}[t]{@{}c@{}}{\normalsize #1}\\[0.3mm]\framebox[9mm]{\rule{0pt}{8mm}}\end{tabular}\hspace{0.6mm}}
\newcommand{\mensajeCifrado}[1]{%
  \par\vspace{2mm}{\centering\fontsize{20}{28}\selectfont\bfseries\addfontfeatures{LetterSpace=12}#1\par}}

% Ficha: un campo por línea.
\newcommand{\campoFicha}[1]{\par\vspace{7mm}\noindent\textbf{#1:}~\hrulefill}

% Compara: dos óvalos que se cruzan -- lo de cada uno a su lado, lo que
% tienen los dos en medio.
\newcommand{\diagramaCompara}[2]{%
  \begin{center}
  \begin{tikzpicture}[line width=1.3pt]
    \draw[colorResponde] (0,0) ellipse (4.3cm and 3.3cm);
    \draw[colorResponde] (4.2,0) ellipse (4.3cm and 3.3cm);
    \node[font=\large\bfseries] at (-1.8,3.7) {#1};
    \node[font=\large\bfseries] at (6,3.7) {#2};
    \node[font=\normalsize, text=colorGris] at (2.1,3.7) {\lblLosDos};
  \end{tikzpicture}
  \end{center}}

% Viñetas: tres recuadros vacíos, cada uno con su pie.
\newcommand{\vineta}[2]{%
  \framebox[\linewidth]{\rule{0pt}{4.8cm}}\par\vspace{1mm}
  {\small\raggedright\textbf{#1.} #2\par}}

% La lupa del bisabuelo relojero -- el símbolo del Club de la Lupa:
% borde dorado, mango de madera oscura con tres rayas grabadas, tal como
% la describe el día 1. Se usa en la portada, el carné del club y el
% diploma. En blanco y negro, el borde y el mango pasan a gris y negro,
% como el resto de la paleta.
\ifbwbook
  \definecolor{colorLupaBorde}{gray}{0.45}
  \definecolor{colorLupaMango}{gray}{0.1}
\else
  \definecolor{colorLupaBorde}{HTML}{C9A227}
  \definecolor{colorLupaMango}{HTML}{5D4037}
\fi
\newcommand{\iconoLupa}[1][1]{%
  \begin{tikzpicture}[scale=#1, line cap=round, line join=round]
    \fill[colorLecturaClaro] (0,0) circle (1.5);
    \draw[colorLupaBorde, line width=#1*5pt] (0,0) circle (1.5);
    \draw[colorLupaMango, line width=#1*13pt] (1.16,-1.16) -- (2.55,-2.55);
    \foreach \d in {0.35,0.65,0.95}
      \draw[colorLupaBorde, line width=#1*1.2pt]
        ({1.33+\d},{-0.99-\d}) -- ({0.99+\d},{-1.33-\d});
    \draw[white, line width=#1*2.5pt] (-0.95,0.55) arc (150:115:1.1);
  \end{tikzpicture}}

\usetikzlibrary{arrows.meta, calc}

% Tamaño del texto del día, según la parte del año (tools/libros.py,
% ENGLISH): de 30-35 palabras en otoño a 60-75 en verano.
\newcommand{\fuenteIngles}[1]{%
  \ifcase#1\relax
  \or \fontsize{20}{27}\selectfont
  \or \fontsize{19}{25.5}\selectfont
  \or \fontsize{18}{24}\selectfont
  \or \fontsize{17}{23}\selectfont
  \fi
}

% Todo lo que la niña lee dentro de una actividad.
\newcommand{\letraActividadIngles}{\fontsize{15}{20}\selectfont}

% La caja de todas las actividades: actividadLupa (preamble-lupa.tex, ver
% allí por qué height fill + height plus) con la letra de este cuaderno.
\newtcolorbox{actividadIngles}[2]{cajabase,
  colback=#1Claro, colframe=#1, colbacktitle=#1, title={#2},
  height fill, height plus=100cm, space to upper, segmentation hidden, middle=1mm,
  fontupper=\letraActividadIngles, fontlower=\letraActividadIngles,
  before upper={\hyphenpenalty=10000\exhyphenpenalty=10000},
  before lower={\hyphenpenalty=10000\exhyphenpenalty=10000}}

% --- piezas de las actividades ---------------------------------------

% La instrucción de cada actividad: en inglés, corta y en negrita -- la
% lee la niña, no el adulto (en "Leo con lupa" iba en pequeño y gris).
\newcommand{\instruccionIng}[1]{{\bfseries #1\par}\vspace{2.5mm}}

% Casilla para marcar, más grande que \casillaMarca (la letra es más
% grande, y la mano que marca, más pequeña).
\newcommand{\casillaIng}{\raisebox{-0.4ex}{\framebox[5.5mm][c]{\rule{0pt}{4.4mm}}}}

% La lista para comprobar un dibujo o para repasar (como listaRepaso de
% preamble.tex), con la casilla grande.
\newenvironment{listaComprueba}{%
  \begin{itemize}[label=\casillaIng, leftmargin=9mm, itemsep=2.5mm, topsep=2mm]
}{%
  \end{itemize}
}

% Sí o no: la frase a la izquierda y las dos casillas a la derecha.
\newcommand{\filaSiNo}[1]{\par\vspace{7mm}\noindent
  \begin{tabularx}{\linewidth}{@{}X r@{}}#1 & \lblSi\ \casillaIng\quad\lblNo\ \casillaIng\end{tabularx}}

% Rodea: las opciones, grandes y separadas, debajo de su pregunta.
\newcommand{\opcionesRodea}[1]{\par\vspace{3mm}{\centering\fontsize{18}{24}\selectfont #1\par}\vspace{7mm}}
\newcommand{\separaOpciones}{\hspace{16mm plus 8mm}}

% Une (relaciona, mitades): un punto al final de cada lado, pegado al
% texto (al final de la última línea a la izquierda, al principio de la
% primera a la derecha), para trazar la línea de punto a punto.
\newcommand{\puntoUne}{\raisebox{0.35ex}{\tikz\fill (0,0) circle (1mm);}}

% Huecos: la caja con las palabras que faltan, y el hueco de cada frase.
\newtcolorbox{cajaPalabrasIng}{enhanced, colback=white, colframe=colorGris,
  boxrule=0.8pt, arc=2mm, halign=center, left=3mm, right=3mm, top=2.5mm, bottom=2.5mm,
  fontupper=\fontsize{17}{22}\selectfont\bfseries, before skip=1mm, after skip=3mm}
\newcommand{\palabrasIng}[1]{\begin{cajaPalabrasIng}#1\end{cajaPalabrasIng}}
\newcommand{\separaPalabras}{\hspace{9mm plus 4mm}}
\newcommand{\huecoIng}{\rule[-0.6mm]{32mm}{0.5pt}}
\newcommand{\separaFrases}{\par\vspace{7mm}}

% Draw the story: como \vineta de preamble-lupa.tex, pero más alta (hay
% más sitio: el texto es más corto) y con el pie a la letra de la
% actividad, que lo lee la niña.
\newcommand{\vinetaIng}[2]{%
  \framebox[\linewidth]{\rule{0pt}{7.5cm}}\par\vspace{1.5mm}
  {\raggedright\textbf{#1.} #2\par}}

% Un dibujo para colorear o una escena donde dibujar (diagrams/english/).
\newcommand{\dibujoIng}[1]{\par\vspace{1mm}{\centering\input{diagrams/english/#1}\par}}

% =========================================================================
% Iconos: uno por actividad, del tamaño de una letra. Sin color propio:
% toman el del texto que los rodea -- blanco en el título de una caja,
% el color de la actividad en la página de las palabras (y gris/negro en
% blanco y negro, como todo lo demás).
% =========================================================================
\tikzset{icono/.style={x=1em, y=1em, line width=0.09em, line cap=round,
  line join=round, baseline=0.05em}}
% Lápiz: dibuja.
\newcommand{\icoLapiz}{\tikz[icono]{%
  \draw (0.08,0.08) -- (0.14,0.36) -- (0.7,0.92) -- (0.92,0.7) -- (0.36,0.14) -- cycle;
  \draw (0.14,0.36) -- (0.36,0.14);
  \fill (0.08,0.08) -- (0.1,0.19) -- (0.19,0.1) -- cycle;}}
% Cera de colores, gruesa y rellena (el lápiz de dibuja es solo el
% contorno): colorea.
\newcommand{\icoCera}{\tikz[icono]{%
  \begin{scope}[rotate around={-45:(0.5,0.5)}]
  \fill (0.33,0.02) rectangle (0.67,0.2);
  \fill (0.33,0.27) rectangle (0.67,0.56);
  \fill (0.33,0.63) -- (0.67,0.63) -- (0.56,0.95) -- (0.44,0.95) -- cycle;
  \end{scope}}}
% Chincheta de mapa: ¿dónde está?
\newcommand{\icoSitio}{\tikz[icono]{%
  \draw (0.5,0.02) .. controls (0.3,0.35) and (0.14,0.5) .. (0.14,0.66)
    arc (180:0:0.36) .. controls (0.86,0.5) and (0.7,0.35) .. (0.5,0.02) -- cycle;
  \fill (0.5,0.66) circle (0.12);}}
% Casilla marcada: sí o no.
\newcommand{\icoMarca}{\tikz[icono]{%
  \draw (0.06,0.06) rectangle (0.94,0.94);
  \draw (0.24,0.5) -- (0.44,0.28) -- (0.78,0.74);}}
% Una letra rodeada: rodea.
\newcommand{\icoRodea}{\tikz[icono]{%
  \draw (0.5,0.45) ellipse (0.55 and 0.42);
  \node[inner sep=0, scale=0.75] at (0.5,0.4) {\bfseries a};}}
% Dos puntos a cada lado, y una línea que ya une dos: relaciona, mitades.
\newcommand{\icoUne}{\tikz[icono]{%
  \fill (0.08,0.8) circle (0.1); \fill (0.08,0.2) circle (0.1);
  \fill (0.92,0.8) circle (0.1); \fill (0.92,0.2) circle (0.1);
  \draw (0.08,0.8) -- (0.92,0.2);}}
% Tres momentos en una flecha: ordena (primero, luego, al final).
\newcommand{\icoOrden}{\tikz[icono]{%
  \draw[-{Latex[length=0.3em, width=0.32em]}] (-0.1,0.5) -- (1.2,0.5);
  \fill (0.12,0.5) circle (0.12); \fill (0.5,0.5) circle (0.12); \fill (0.88,0.5) circle (0.12);}}
% Una palabra a la que le falta un trozo: huecos.
\newcommand{\icoHueco}{\tikz[icono]{%
  \node[inner sep=0, anchor=base] at (0.05,0.1) {a};
  \draw (0.32,0.08) -- (0.78,0.08);
  \node[inner sep=0, anchor=base] at (1.05,0.1) {c};}}
% Lupa: busca palabras en el texto.
\newcommand{\icoBusca}{\tikz[icono]{%
  \draw (0.42,0.58) circle (0.32);
  \draw[line width=0.16em] (0.64,0.36) -- (0.92,0.08);}}
% Interrogación: adivina.
\newcommand{\icoAdivina}{\tikz[icono]{%
  \draw (0.5,0.5) circle (0.46);
  \draw (0.34,0.64) arc (160:-60:0.18) -- (0.5,0.38);
  \fill (0.5,0.22) circle (0.06);}}
% Tres viñetas: dibuja la historia.
\newcommand{\icoVinetas}{\tikz[icono]{%
  \draw (0.02,0.2) rectangle (0.3,0.8);
  \draw (0.36,0.2) rectangle (0.64,0.8);
  \draw (0.7,0.2) rectangle (0.98,0.8);}}
% Estrella: mira atrás, repasa.
\newcommand{\icoEstrella}{\tikz[icono]{%
  \fill (0.5,0.97) -- (0.62,0.63) -- (0.97,0.63) -- (0.69,0.42) -- (0.8,0.07)
    -- (0.5,0.28) -- (0.2,0.07) -- (0.31,0.42) -- (0.03,0.63) -- (0.38,0.63) -- cycle;}}
% Cara sonriente: tu turno.
\newcommand{\icoCara}{\tikz[icono]{%
  \draw (0.5,0.5) circle (0.46);
  \fill (0.35,0.62) circle (0.06); \fill (0.65,0.62) circle (0.06);
  \draw (0.28,0.4) arc (200:340:0.24);}}
% Libro abierto: lee (la caja verde del texto, en la página de las
% palabras) -- el mismo icono de la portada del cuaderno de frases.
\newcommand{\icoLibro}{\tikz[icono]{%
  \draw (0.02,0.85) -- (0.5,0.95) -- (0.98,0.85) -- (0.98,0.12) -- (0.5,0.22) -- (0.02,0.12) -- cycle;
  \draw (0.5,0.95) -- (0.5,0.22);}}
